
%(BEGIN_QUESTION)
% Copyright 2006, Tony R. Kuphaldt, released under the Creative Commons Attribution License (v 1.0)
% This means you may do almost anything with this work of mine, so long as you give me proper credit

En operatør som kjeder seg fyller en stor tank med vann fra ulike kilder, der strømningsratene fra kildene varierer over tid. Han fordriver tiden med å skrive ned volumverdier fra nivåindikatoren (kalibrert i gallon) ved ulike tidspunkter, og noterer klokkeslett ved siden av volumene:

% No blank lines allowed between lines of an \halign structure!
% I use comments (%) instead, so that TeX doesn't choke.

$$\vbox{\offinterlineskip
\halign{\strut
\vrule \quad\hfil # \ \hfil & 
\vrule \quad\hfil # \ \hfil \vrule \cr
\noalign{\hrule}
%
% First row
Nivåindikator & Tid \cr
% 
(gallons) & (time:minutt) \cr
%
\noalign{\hrule}
%
% Another row
120.7 & 9:17 \cr
%
\noalign{\hrule}
%
% Another row
1005.4 & 9:21 \cr
%
\noalign{\hrule}
%
% Another row
1377.8 & 9:23 \cr
%
\noalign{\hrule}
%
% Another row
2050.2 & 9:26 \cr
%
\noalign{\hrule}
%
% Another row
2944.6 & 9:30 \cr
%
\noalign{\hrule}
%
% Another row
4875.1 & 9:40 \cr
%
\noalign{\hrule}
%
% Another row
5101.8 & 9:45 \cr
%
\noalign{\hrule}
} % End of \halign 
}$$ % End of \vbox

Beregn gjennomsnittlig vannstrøm inn i tanken mellom følgende tidspunkter:

\begin{itemize}
\item{} Mellom 9:17 og 9:21, gjennomsnittsstrøm = \underbar{\hskip 50pt} GPM
\vskip 5pt
\item{} Mellom 9:21 og 9:23, gjennomsnittsstrøm = \underbar{\hskip 50pt} GPM
\vskip 5pt
\item{} Mellom 9:23 og 9:26, gjennomsnittsstrøm = \underbar{\hskip 50pt} GPM
\vskip 10pt
\item{} Mellom 9:17 og 9:26, gjennomsnittsstrøm = \underbar{\hskip 50pt} GPM
\end{itemize}

Sammenlign deretter gjennomsnittsstrømmen i de tre første intervallene med gjennomsnittsstrømmen over summen av disse intervallene (9:17 til 9:26). Hva forteller dette oss om beregning av vannstrøm basert på volum- og tidsmålinger?

\vskip 20pt \vbox{\hrule \hbox{\strut \vrule{} {\bf Forslag til sokratisk diskusjon} \vrule} \hrule}

\begin{itemize}
\item{} Identifiser de aritmetiske operasjonene som trengs for å beregne endringsrater, slik som strømning.
\item{} Denne typen repeterende beregning egner seg godt for programmerbar kalkulator eller et {\it regneark} på PC. Hvis du kan litt regneark, prøv å lage ett som beregner gjennomsnittlig strømningsrate fra denne volumtabellen.
\end{itemize}

\underbar{file i01503}
%(END_QUESTION)





%(BEGIN_ANSWER)

\begin{itemize}
\item{} Mellom 9:17 og 9:21, gjennomsnittsstrøm = \underbar{\bf 221.18} GPM
\vskip 5pt
\item{} Mellom 9:21 og 9:23, gjennomsnittsstrøm = \underbar{\bf 186.2} GPM
\vskip 5pt
\item{} Mellom 9:23 og 9:26, gjennomsnittsstrøm = \underbar{\bf 224.13} GPM
\vskip 10pt
\item{} Mellom 9:17 og 9:26, gjennomsnittsstrøm = \underbar{\bf 214.39} GPM
\end{itemize}

%(END_ANSWER)





%(BEGIN_NOTES)

For å beregne hver gjennomsnittsstrømningsrate må vi dele differansen i volum ($\Delta V$) på tilsvarende differanse i tid ($\Delta t$):

$$\overline{Q} = {\Delta V \over \Delta t}$$

\vskip 10pt

Lærdommen er at vi får mer detaljert informasjon om strømning hvis vi måler volum (og tid) {\it oftere}. Målinger med lange mellomrom viser bare et {\it gjennomsnitt} over tid, og vi mister det ``fine'' bildet av hvordan strømningen topper og avtar. Den ytterste konsekvensen er å redusere tiden mellom samplingsintervaller til vi {\it kontinuerlig} måler differensialer i volum ($dV$) og tid ($dt$) for å få ekte strømningsberegning:

$$Q = {dV \over dt}$$


%INDEX% Mathematics, calculus: derivative (calculating flow rates from measured volumes at specific times)

%(END_NOTES)

